\section{Extension: how well are the crop parameters supported?}\label{sec:parevidence}
\paragraph{Purpose.} The index and the future yield accounting rest on crop parameters of very different quality.
Some are printed in a peer-reviewed experiment on the same crop, outcome and temperature metric. Others are our
conversions, values borrowed from another crop, or assumptions. This section audits every parameter the model uses,
states how well each is supported, and shows what the headline result looks like when parameters without defensible
support are left out instead of being filled in.

\subsection{Method}
\paragraph{Register.} The register has \parNRows\ rows: one per crop and parameter actually used (heat threshold
$T_{\text{crit}}$, heat slope $s_T$, $K_y$, $K_c$, stage days, calendar, flowering window, temperature-metric
conversion, food-item scope), plus one row per perturbation range we have run. Each row records the value and units,
the source, the status, whether the primary source was read and at what level, the experimental setting, how well it
applies to our use, any derivation, limitations, candidate replacements and the final treatment. The register is
\file{docs/evidence/parameter\_evidence\_register.csv}; the reasoning and quotations are in
\file{docs/evidence/parameter\_evidence\_audit.md}.

\paragraph{Definitions.}
\begin{itemize}
\item \textbf{Published:} the value is printed in a source for the same species, outcome and metric.
\item \textbf{Derived:} our arithmetic or fit on published data (for example the conversion of a chamber threshold to
$T_{\text{eff}}$; refitted by \file{code/25\_parameter\_fits.py}).
\item \textbf{Borrowed:} another crop's, outcome's or setting's value.
\item \textbf{Assumption:} our own choice, such as a diurnal range.
\item \textbf{Unsupported:} no source supports it for this use.
\end{itemize}
Applicability (good, partial, poor) compares the source with our use on species and cultivar, measured outcome,
temperature or water metric, growth stage and exposure duration, experimental setting, and units.

\begin{table}[H]\centering\small
\caption[Register by status and applicability]{The \parNRows\ register rows by status and applicability. Rows
without an applicability grade are perturbation ranges.}\label{tab:parstatus}
\adjustbox{max width=\linewidth}{\begin{tabular}{@{}lrrrrr@{}}
\toprule
Status & Rows & Good & Partial & Poor & Not assessed \\
\midrule
Published & 38 & 27 & 8 & 2 & 1 \\
Derived & 18 & 5 & 13 & 0 & 0 \\
Borrowed & 13 & 0 & 10 & 1 & 2 \\
Assumption & 27 & 2 & 7 & 2 & 16 \\
Unsupported & 1 & 0 & 0 & 1 & 0 \\
All & 97 & 34 & 38 & 6 & 19 \\
\bottomrule
\end{tabular}
}
\end{table}

\subsection{What the audit found}
Table~\ref{tab:parkey} lists the items the project flagged as most uncertain. Quotations, page numbers and read levels
are in the audit file.

\begin{table}[H]\centering\scriptsize
\caption[Audit of the most uncertain parameters]{Audit of the most uncertain parameters. ``Treatment'' is the final
treatment in this report.}\label{tab:parkey}
\begin{tabular}{@{}p{2.2cm}p{1.9cm}p{6.7cm}p{3.9cm}@{}}
\toprule
Parameter & Current value & Finding & Treatment \\
\midrule
Rice $K_y$ & 1.25 (maize's) & FAO gives no rice $K_y$; FAO-33 only classes rice as ``high''. No field seasonal $K_y$ on the FAO definition was found. The only rice values found are stage-wise (1.14, 1.07, 1.02) and partly model-derived \citep{anupoju2020}. & Borrowed proxy, kept and labelled. The source text in the ledger was wrong (it implied an FAO rice value) and is corrected. \\
Rice $s_T$ & 0.07 per $^\circ$C above 35\,$^\circ$C $T_{\text{eff}}$ & A spikelet-fertility slope of one genotype (IR64; tissue temperature above 29.6\,$^\circ$C, hours at anthesis; \citealp{jagadish2007}). Fertility is one yield component and does not translate one-to-one into yield \citep{prasad2006rice,hatfield2011}. & \textbf{Unsupported as a yield slope on $T_{\text{eff}}$.} Heat component unavailable in the screened view. \\
Rice yield per $^\circ$C & (none) & Season-long sunlit-chamber work gives about $-10\%$ per $^\circ$C of mean temperature above $25\,^\circ$C, but it was read only second-hand, in \citet{hatfield2011}, and the chamber setting conflicts with field evidence on maximum temperature \citep{welch2010}. & Candidate recorded, \emph{not} adopted and not computed. \\
Rapeseed $s_T$ & 0.10 (wheat's span) & No \emph{Brassica} yield slope found in the papers checked \citep{morrison2002,kutcher2010,gan2004,angadi2000,nuttall1992}, and no Indian mustard dose--response. & \textbf{Unsupported.} Heat component unavailable in the screened view. \\
Rapeseed $T_{\text{crit}}$ & 26.5\,$^\circ$C $T_{\text{eff}}$ & Converted from 29.5\,$^\circ$C daily maximum \citep{morrison2002} with an unsourced 12\,$^\circ$C diurnal range. The experiment included \emph{B.~juncea} but Canadian cultivars only. & Assumption; partly applicable to Indian mustard; kept as a labelled proxy. \\
Conversion to $T_{\text{eff}}$ & as derived & Exact algebra for 12\,h/12\,h step regimes, resting on assumptions: constant chamber regimes stand for field monthly $T_{\text{eff}}$, and a daily accumulation threshold is reused as a monthly-mean threshold. Not a published relation. & Derived calculation, labelled; not rice (rice was never converted). \\
Lentil $K_y$ & 1.0 & A field study \citep{nandi2021}; no unstressed control; treatments mix tillage and fertiliser. Partial match to the FAO definition. No other lentil $K_y$ found. & Partial; proxy. \\
Chamber to field & legume thresholds, slopes & The direction is supported: field thresholds sit 2--4\,$^\circ$C above chamber ones \citep{devasirvatham2013,morrison2002}. The $\times0.7$ slope factor has no source. & Exploratory perturbation only. \\
Bean in India and Myanmar & \emph{Phaseolus} parameters & FAO's item definition says several countries include \emph{Vigna} beans \citep{fao1994commodities}. India and Myanmar are ranks 1 and 3 of the bean panel. An unreviewed USDA statement puts \emph{Vigna} at about 93\% of India's figure \citep{phaseoluscgc2020}. & \textbf{Pairs flagged not applicable} in screened view V2. \\
Calendars & cool-season rule & No gridded observed calendar exists for chickpea or lentil. A GGCMI Phase 3 rapeseed calendar exists \citep{ggcmi3calendar} but gives maturity, not harvest, and a Canadian planting day (91) far from the rule's. & Proxy kept, labelled; alternative not adopted. \\
\bottomrule
\end{tabular}
\end{table}

\paragraph{Ranges are perturbations, not uncertainty.} Of the perturbation rows, \parNExploratory\ are exploratory:
a chosen offset or factor with no literature support as a range. Only \parNLitRange\ rows carry a literature value
(bean $K_y$ 0.6--1.4 from two field values, and treatment values of chickpea and lentil $K_y$), and none of these is
a confidence interval. The earlier statement that wheat stays first under every robustness test therefore remains true, but it is
a test of sensitivity to our choices, not an uncertainty analysis. The same holds for the Monte Carlo over index
weightings: its distributions (uniform level-versus-change and heat-versus-water mixes, Dirichlet nutrient weights)
are ours, so its frequencies describe sensitivity to our design, not probabilities. No new distribution or Monte
Carlo run was added in this extension.

\paragraph{What is at stake.} The two unsupported slopes matter very differently for the present and the future.
\begin{itemize}
\item \textbf{Present scores:} heat is only \parHeatPresRice\% of rice's and \parHeatPresRape\% of rapeseed's F4,
because their thresholds are rarely exceeded in monthly means. Their present scores are essentially water scores and
are barely affected by the unsupported slopes.
\item \textbf{Projected change:} heat is \parHeatFutRape\% of rapeseed's change but only \parHeatFutRice\% of rice's
(shares of the sum of absolute pair changes). The statement that rapeseed's projected increase is heat-driven rests
on the unsupported slope and cannot be treated as a finding.
\item \textbf{Everything:} common bean sets the heat reference for all crops (Section~\ref{sec:exposure}), and bean in
India and Myanmar is the least applicable part of the bean evidence.
\end{itemize}

\subsection{Evidence-screened view}
Where no defensible replacement exists, the gap is recorded and the affected result is left unavailable. Two rules
are applied to the register, without filling or re-weighting:
\begin{itemize}
\item \textbf{Rule A:} a hazard component is unavailable for a crop when the register marks a parameter it needs as
unsupported (rice and rapeseed heat slope). The pair's F4 is then unavailable. The water term is not up-weighted
and the missing heat term is not set to zero, because ``not estimated'' is not ``no effect''.
\item \textbf{Rule B:} a pair is flagged not applicable when its crop definition does not match the parameter source
(bean in India and Myanmar).
\end{itemize}
Three views are reported, none of them a nine-crop ranking except the first.
\begin{itemize}
\item \textbf{V0:} the full index as in Section~\ref{sec:present}. It is an \emph{assumption-dependent} (legacy)
result: it uses proxies and borrowed values.
\item \textbf{V1:} rule A applied. Rice and rapeseed are unavailable, so only seven crops are compared.
\item \textbf{V2:} rules A and B applied, so bean is computed over its remaining five countries.
\end{itemize}
Reference maxima are recomputed over the pairs that remain, which changes the normalisation and is documented as
decision E20. In V1 the heat, water and sensitivity maxima are unchanged, but the importance maximum falls because
rice in Bangladesh, a removed pair, set it; every F4 is rescaled by the common factor \scrRescale\ (checked), so the
order of the remaining crops cannot change. In V2 the heat maximum falls from \scrVZeroHeatMax\ to
\scrVTwoHeatMax\,$^\circ$C (it is now set by \scrVTwoHeatPair) because bean in India no longer sets it, which makes
every exposure to heat \scrHeatRatio\ times larger; wheat in India rises from \scrWheatIndVzero\ to
\scrWheatIndVtwo.

\begin{table}[H]\centering\scriptsize
\caption[Evidence-screened F4]{Crop scores (rank in brackets among the crops with a score) in the full index (V0),
without rice and rapeseed (V1), and also without bean in India and Myanmar (V2). Equal-weighted and
production-weighted. ``Unavailable'' is not ``lowest''. Last column: pairs with a score in V2 out of all modelled
pairs of the crop.}\label{tab:screened}
\adjustbox{max width=\linewidth}{\begin{tabular}{@{}lrrrrrrr@{}}
\toprule
Crop & V0 eq. & V1 eq. & V2 eq. & V0 prod. & V1 prod. & V2 prod. & V2 pairs \\
\midrule
Wheat & 0.0563 (1) & 0.0689 (1) & 0.0760 (1) & 0.0696 (1) & 0.0852 (1) & 0.0946 (1) & 6/6 \\
Rice & 0.0024 (7) & unavailable & unavailable & 0.0015 (9) & unavailable & unavailable & 0/7 \\
Maize & 0.0297 (2) & 0.0363 (2) & 0.0363 (2) & 0.0080 (6) & 0.0098 (5) & 0.0098 (5) & 6/6 \\
Chickpea & 0.0053 (6) & 0.0064 (5) & 0.0064 (5) & 0.0228 (3) & 0.0279 (3) & 0.0279 (3) & 7/7 \\
Lentil & 0.0021 (8) & 0.0025 (6) & 0.0025 (6) & 0.0040 (7) & 0.0048 (6) & 0.0048 (6) & 7/7 \\
Common bean & 0.0172 (3) & 0.0210 (3) & 0.0293 (3) & 0.0233 (2) & 0.0285 (2) & 0.0466 (2) & 5/7 \\
Soybean & 0.0078 (4) & 0.0096 (4) & 0.0096 (4) & 0.0100 (5) & 0.0122 (4) & 0.0122 (4) & 7/7 \\
Groundnut & 0.0015 (9) & 0.0018 (7) & 0.0020 (7) & 0.0018 (8) & 0.0022 (7) & 0.0025 (7) & 5/5 \\
Rapeseed-mustard & 0.0078 (5) & unavailable & unavailable & 0.0157 (4) & unavailable & unavailable & 0/7 \\
\bottomrule
\end{tabular}
}
\end{table}

\paragraph{Result.} Wheat ranks first in all three views under both country weightings, and the order of the other
crops that remain does not change between V1 and V2. This is a statement about the seven crops whose parameters
survive the screen. It says nothing about rice or rapeseed, whose position cannot be established from the
supported evidence, and it does not turn V1 or V2 into a nine-crop comparison.

\paragraph{Remaining gaps.} Unresolved, and left open rather than filled:
\begin{itemize}
\item a supported yield slope and threshold for rice on a metric convertible to $T_{\text{eff}}$;
\item any heat slope for rapeseed and Indian mustard;
\item a \emph{Vigna} parameter set for the bean pairs in India and Myanmar (not assessed);
\item observed calendars for chickpea and lentil;
\item full texts of \citet{prasad2002bean}, \citet{prasad2003peanut} and \citet{morrison2002}, which remain closed
access and are taken from abstracts and from \citet{boote2005}.
\end{itemize}
A rice sensitivity run on the Baker-derived candidate would need a new hazard run and rests on a secondary source;
it was not made.
